% Working mathematical development. Epsilon convention: 1 iff equality.
\section{Scalar self-seeded iteration with actual equality}
\label{sec:selfseed-theory}
There is one numerical sort, with terms
\[
 t ::= x\mid 0\mid St\mid \epsilon(t,t)\mid F(n;v.h).
\]
Only $h$ binds $v$. Every fixed earlier-stage body $h(v,\vec p)$ gives
one first-order symbol $F_h(n,\vec p)$; there is no function-variable
sort and no uniformly total evaluator. Ordinary intuitionistic
first-order logic with actual equality is used.

The numerical axioms are
\begin{align*}
F(0;v.h)&=S(h[0/v]),&F(Sn;v.h)&=h[F(n;v.h)/v],\\
Sx=Sy&\longrightarrow x=y,&0=S0&\longrightarrow\bot,\\
\epsilon(x,y)=0&\ \lor\ \epsilon(x,y)=1,&
\epsilon(x,y)=1&\longleftrightarrow x=y.
\end{align*}
Here $1=S0$ and all free variables are universally closed.
The sole induction rule is contextual induction on a literal equation:
\[
\frac{\Gamma\vdash E(0)\quad\Gamma,E(z)\vdash E(Sz)}
     {\Gamma\vdash E(t)},
\]
where $z$ is fresh for the retained context and fixed parameters.
All substitutions avoid capture. Successor injectivity is an explicitly
charged additional axiom. Numerical coverage is derived below rather
than postulated.

\paragraph{Optional definitional removal of the successor constructor.}
The source grammar may be shortened to $x,0,\epsilon,F$ by defining
$S(a):=F(0;v.a)$. The constant-body computation already proves this
identity in the displayed theory. In the shortened presentation write
both computation schemas and the boundary axioms using that abbreviation.
The fixed constant-body symbol belongs to the first description stage;
its own base equation is tautological, while its step equation and the
boundary axioms remain substantive. Reintroducing $S$ by the definition
recovers the displayed theory, and expanding $S$ gives the converse
translation. This familiar constructor-elimination device already occurs
in the supplied successor-seeded drafts; the new contribution below is
the scalar self-seeded completion, not this abbreviation by itself.

\subsection{Elementary Boolean operations and quantifier-free induction}
Put
\[
D(n,a):=F(n;v.a),\qquad Z(n):=D(n,0).
\]
The computation axioms prove
\[
D(0,a)=Sa,\quad D(Sn,a)=a,\quad Z(0)=1,\quad Z(Sn)=0.
\]
If $Sn=0$, then $D(Sn,0)=0=D(0,0)=1$, contradicting separation.
Thus every successor is nonzero. Injectivity and this consequence
distinguish every fixed pair of distinct numerals.

For Boolean $b,c$ define
\[
\neg_B b=\epsilon(b,0),\qquad
b\mathbin{\wedge_B}c=\epsilon(D(b,\epsilon(c,0)),0).
\]
Their four truth-table cases follow from the displayed computations.
Disjunction is defined by De Morgan on these Boolean values, and
implication by $\neg_Bb\vee_Bc$. Numerical coverage is not substituted
for Boolean range: the primitive range axiom supplies the two cases
of every equality characteristic, and the finite truth tables preserve
that range.

For each fixed quantifier-free formula $A$, an external recursion on
syntax produces a Boolean source term $\chi_A$ such that
\[
 \chi_A=1\longleftrightarrow A,\qquad
 \chi_A=0\longleftrightarrow\neg A.
\]
The implication clause is intuitionistically valid because its
quantifier-free antecedent has the proved Boolean characteristic and
is therefore decidable. Equation induction on $\chi_A=1$ gives
contextual quantifier-free induction. This does not add induction for
arbitrary quantified formulas.

\subsection{A finite-case predecessor graph}
For a passive parameter $a$ define
\[
 U_a(v):=D(\epsilon(Sv,a),v),\qquad
 V(j,a):=F(j;v.U_a(v)),\qquad W(a):=V(a,a).
\]
If $a\ne1$, then $V(0,a)=2$. The Boolean cases show that $U_a(v)$
is $v$ when $Sv=a$, and $Sv$ otherwise. Contextual quantifier-free
induction on $j$ therefore proves
\[
a\ne1\quad\longrightarrow\quad
\bigl(V(j,a)=S^2j\ \lor\ S(V(j,a))=a\bigr).
\tag{P-invariant}
\]
In the first alternative of the step, a successful test gives the
second alternative and a failed test advances the clock by one. In
the second alternative the update fixes the value.

Take $a=S^3m$ and $j=m$. This $a$ differs from $1$ by injectivity and
successor nonzeroness. The first alternative gives $V(m,a)=S^2m$.
The second gives the same equation by successor injectivity, since
$S(V(m,a))=S^3m$. The update now fixes the value, so exactly three
further computation steps give
\[
                         W(S^3m)=S^2m.\tag{W}
\]
No induction with a changing passive parameter was used: the induction
proved P-invariant with $a$ fixed, and only afterwards specialized it.

The values at the three small arguments, calculated directly, are
$W(0)=2$, $W(1)=2$, and $W(2)=4$. Define the graph $\Phi_P(n,p)$ by
the following finite quantifier-free cases:
\[
\begin{split}
\Phi_P(n,p):={}&(n=0\wedge p=0)\vee(n=1\wedge p=0)
 \vee(n=2\wedge p=1)\\
 &\vee(n\ne0\wedge n\ne1\wedge n\ne2\wedge p=W(n)).
\end{split}
\]
Equality decidability makes these guards exhaustive; they are pairwise
disjoint. Thus the graph is provably total and single-valued, without
an existence oracle or a new recursion schema. Introduce its function
$P$ in a conservative definitional extension. No constructor-coverage
principle has yet been used.

Every term containing finitely many $P$ occurrences can be expanded
into finitely many quantifier-free source guards with source value
terms: refine argument cases and then the four cases of $\Phi_P$.
Consequently each quantifier-free formula of the extension has a
provably equivalent source quantifier-free formula. Applying source
quantifier-free induction proves its extension instance. This is a
finite-case argument, not the invalid assertion that every conservative
definition automatically preserves a restricted induction schema.

\paragraph{Deriving predecessor and coverage together.}
Apply this legitimate quantifier-free induction in the extension to
the joint invariant
\[
 \bigl(n=0\ \lor\ n=S(Pn)\bigr)\quad\wedge\quad P(Sn)=n.
 \tag{Joint invariant}
\]
At zero the first conjunct is immediate, and the second is $P1=0$
from the finite cases. For the successor first conjunct, the old
second conjunct gives $Sn=S(P(Sn))$. For its second conjunct
$P(S^2n)=Sn$, split the old first conjunct. If $n=0$, use $P2=1$.
If $n=S(Pn)$, then $S^2n=S^3(Pn)$ is distinct from $0,1,2$;
the last predecessor case and (W) give
\[
 P(S^2n)=W(S^3(Pn))=S^2(Pn)=Sn.
\]
Thus the source proves, through finite-case elimination,
\[
 P0=0,\qquad P(Sn)=n,\qquad n=0\ \lor\ n=S(Pn).
\]
Ordinary existential introduction now gives numerical coverage
$n=0\lor\exists m\,n=Sm$. This is an axiom deletion, not a
standard-model observation: the single displayed joint induction
constructs both conclusions internally, before using coverage in
any subsequent order argument.

\subsection{Capped counting and an order candidate}
Define
\[
H_a(v):=D(\epsilon(v,S^2a),v),\qquad C(n,a):=F(n;v.H_a(v)).
\]
Since $S^2a\ne0$, $H_a(0)=1$ and $C(0,a)=2$. The same two-case
induction as above gives
\[
C(n,a)=S^2n\ \lor\ C(n,a)=S^2a,\qquad C(a,a)=S^2a.
\tag{C-invariant}
\]
If the cap has been reached the next state is identical. In particular
$C(n,0)=2$ by equation induction.

Coverage proves $D(c,Sv)=S(D(c,v))$. Boolean reflection and successor
injectivity prove
$\epsilon(Sv,S^3a)=\epsilon(v,S^2a)$; hence
$H_{Sa}(Sv)=S(H_a(v))$. Moreover $C(1,Sa)=3$ because the cap
$S^3a$ cannot equal $2$. Equation induction on $n$, keeping $a$ fixed,
now proves the factor identity
\[
                            C(Sn,Sa)=S(C(n,a)).\tag{C-factor}
\]
Thus the Boolean term
\[
                          L(a,n):=\epsilon(C(n,a),S^2a)
\]
satisfies the familiar comparison recurrences
\[
 L(0,n)=1,\qquad L(Sa,0)=0,\qquad L(Sa,Sn)=L(a,n).
\tag{L-recurrence}
\]
Write $a\le n$ for $L(a,n)=1$. The orbit supplies more than its
two-variable comparison recurrences. Its persistence and the invariant
give the unilateral law
\[
 a\le Sb\quad\longleftrightarrow\quad (a\le b\ \lor\ a=Sb).
 \tag{New-cap}
\]
Indeed, if the cap is already reached it persists. If it is not reached,
the old value is $S^2b$, the next value is $S^3b$, and this equals the
cap $S^2a$ exactly when $a=Sb$. The reverse direction uses persistence
and $C(a,a)=S^2a$.

Linearity now has a valid fixed-parameter induction. Fix $a$ and induct
on $n$ in $a\le n\lor n\le a$. The base uses $0\le a$. At the
successor, $a\le n$ persists. In the other alternative split $a=0$
or $a=Sb$ by coverage. The zero case gives $n=0$. In the successor
case New-cap gives $n\le b$ or $n=a$. The first gives $Sn\le Sb=a$
by the factor identity; the second gives $a\le Sn$ by reflexivity and
persistence. The induction hypothesis has kept $a$ fixed throughout.

For completeness, $Sn\not\le n$ follows by equation induction from
$L(1,0)=0$ and the factor recurrence. The implication
$Sn\le a\Rightarrow n\le a$ follows by splitting $a=0$ or $a=Sb$,
using the factor recurrence and persistence. Antisymmetry is induction
on $n$, with $a$ fixed, in
$(a\le n\wedge n\le a)\to a=n$. At a successor New-cap splits
$a\le Sn$ into $a=Sn$ or $a\le n$; in the latter case $Sn\le a$
implies $n\le a$, the induction hypothesis gives $a=n$, and
$Sn\not\le n$ gives a contradiction. Transitivity is induction on
$c$, with $a,b$ fixed, in
$(a\le b\wedge b\le c)\to a\le c$. New-cap splits the new second
comparison into its old one or $b=Sc$. These give the conclusion by
the induction hypothesis and persistence, or by substitution.
Thus $\le$ is an internal discrete linear order with least element
zero. Successor preserves and reflects it by L-recurrence.

The finite-case graph choosing $n$ when $n\le a$ and $a$ otherwise
defines $\min(n,a)$. C-invariant, antisymmetry, and the definition
of $L$ prove
\[
                          C(n,a)=S^2\min(n,a).
\tag{C-min}
\]
All subsequent uses of order and minimum refer to these derived
graphs and laws, not standard-model interpretations.

\subsection{Conditional bounded selection: exact local lemma}
Let $M(c,p,q)=(c\wedge_Bp)\vee_B(\neg_Bc\wedge_Bq)$, for Boolean
$c,p,q$. Define the following source terms:
\[
\begin{split}
T(b,c,a,d)&:=F\bigl(b;v.D(M(c,\epsilon(v,S^2a),
                                  \epsilon(v,S^2d)),v)\bigr),\\
\operatorname{Sel}(b,c,a,d)&:=W(W(T(b,c,a,d))).
\end{split}
\]
If $c=1$, recurrence uniqueness gives $T(b,c,a,d)=C(b,a)$;
if $c=0$, it gives $T(b,c,a,d)=C(b,d)$. Accordingly the exact proved
conditional hypotheses
\[
 a=S^2a_0,\ d=S^2d_0,\ L(a,b)=1,\ L(d,b)=1
\]
imply that Sel returns $a$ when $c=1$ and $d$ when $c=0$.
Indeed the selected capped value is $S^2a$ or $S^2d$, and both
applications of $W$ are in the range covered by (W), so they remove
exactly two successors. This lemma uses the literal equations
$L(a,b)=L(d,b)=1$, not unproved order properties.

The stronger branch-local version is also valid: when $c=1$ only
$2\le a\le b$ is needed, and when $c=0$ only $2\le d\le b$ is
needed. The Boolean mux erases the other test pointwise; recurrence
uniqueness then identifies $T$ with the selected capped counter. The
compiler below nevertheless makes both candidates bounded before
using Sel, to keep its invariant uniform.

\section{A source compiler for bounded scalar iteration}
\label{sec:selfseed-bounded}
This section is the step that allows graph-defined operations to be
used in updates without enlarging the primitive recursor schema.

\subsection{Finite cases and encoded values}
Call a total single-valued function graph \emph{finite-case} if it has
finitely many source quantifier-free guards $A_i(\vec x)$ with source
value terms $t_i(\vec x)$, the guards are provably exhaustive, and
the values agree whenever their guards overlap. Equivalently its
graph is $\bigvee_i(A_i(\vec x)\wedge y=t_i(\vec x))$.
Finite composition of such graphs again has this form, by taking
finite common refinements of the argument cases. The same expansion
eliminates the corresponding graph-defined function symbols from
every fixed quantifier-free formula. Hence quantifier-free induction
transfers to their conservative definitional extension. This assertion
concerns ordinary function terms only: it does not freely admit a
new primitive $F$ body containing an opaque defined function.

Write $e(x)=S^2x$. On encoded inputs $u=e(x)$ the following finite
source cases decode $x$:
\[
 u=2:\ 0;\qquad u=3:\ 1;\qquad u\ne2,3:\ W(W(u)).
 \tag{Decode}
\]
The last case has $u\ge4$ by coverage, and two applications of (W)
therefore remove exactly two successors. No source term is claimed
to be a correct decoder at all encoded values without these cases.

For $B\ge2$ define the source clipping term
\[
                    K_B(v):=W(W(C(v,B))).
\]
When $v\ge2$, C-min and (W) prove
\[
             K_B(v)=\min(v,B),\qquad 2\le K_B(v)\le B.
 \tag{Clip}
\]
Only arguments of the form $e(t)$ or values already in $[2,B]$
are supplied to this term in the following compiler. There is no
assertion that Clip has the same description at $v=0,1$.

Let $g(\vec p)$ and $h(v,\vec p)$ be fixed finite-case graphs. Fix a
raw cap $b$ and put $B=e(b)$. Expand the finite cases of $g$, encode
and clip each candidate value directly as $C(t_i,b)=e(\min(t_i,b))$, and fold the
cases by nested $\operatorname{Sel}(B,\chi_{A_i},-, -)$, with
default value $2$. This gives a source term $a_B(\vec p)$ in $[2,B]$
that equals $e(\min(g(\vec p),b))$.

For the update, first compose the finite cases of $h$ with Decode
for its state argument, refine the resulting source guards, and
perform the same encoded clipping and selection. The result is a
literal original-language source term $h_B(u,\vec p)$, provably
in $[2,B]$ for every numerical $u$, such that
\[
 h_B(e(v),\vec p)=e(\min(h(v,\vec p),b)).
 \tag{Compiled update}
\]
The universal range follows because every candidate is $C(t_i,b)$
and hence lies in $[2,B]$, including at inputs not satisfying its
guard. Equality characteristics and Boolean connectives make every
guard test Boolean; no lazy-evaluation assumption is involved.

\subsection{Two transient sentinels}
Define the source term
\[
\begin{split}
 H_B(u):=\operatorname{Sel}\bigl(S^2B,\epsilon(u,0),B,\quad\qquad\qquad\qquad\qquad\qquad\\
 \operatorname{Sel}\bigl(S^2B,\epsilon(u,SB),S^2B,
 \operatorname{Sel}(S^2B,\epsilon(u,S^2B),a_B,h_B(u))\bigr)\bigr).
\end{split}
\]
Every candidate of these selectors lies in $[2,S^2B]$. Since the
sentinels $SB,S^2B$ are distinct and exceed $B$, the theory proves
\[
 H_B(0)=B,\quad H_B(SB)=S^2B,\quad H_B(S^2B)=a_B,
 \quad 2\le u\le B\Longrightarrow H_B(u)=h_B(u).
\]
Consequently
\[
 F(0;u.H_B)=SB,\quad F(1;u.H_B)=S^2B,\quad F(2;u.H_B)=a_B.
\]
Set
\[
                      Y_b(n,\vec p):=F(e(n);u.H_B(u)).
\]
Quantifier-free induction proves $2\le Y_b(n,\vec p)\le B$.
Its base and step equations are
\[
 Y_b(0,\vec p)=a_B,\qquad
 Y_b(Sn,\vec p)=h_B(Y_b(n,\vec p),\vec p).
\]
Applying the finite Decode cases to $Y_b$ produces a finite-case
graph $R_b(n,\vec p)$ satisfying
\[
\begin{split}
 R_b(0,\vec p)&=\min(g(\vec p),b),\\
 R_b(Sn,\vec p)&=\min(h(R_b(n,\vec p),\vec p),b),\\
 R_b(n,\vec p)&\le b.
\end{split}
\tag{Bounded iteration}
\]
Every occurrence of $F$ in $Y_b$ has a literal source body. In
particular, the finitely many auxiliary graph functions do not occur
as unexplained new symbols inside that body. The construction is an
external finite compiler from their cases to source syntax.

\paragraph{Scope of the lemma.}
This gives arbitrary initialized, capped \emph{scalar} iteration of
finite-case graph functions. It does not yet assert unbounded
recursion totality. Below each unbounded use is assigned a source
unary cap and supplied with an internal proof that clipping cannot
occur on its intended arguments.

\section{Arithmetic graphs without a circular size bound}
\label{sec:selfseed-arithmetic}

All graph-defined functions in this section are finite-case, and every
induction mentioned is quantifier-free after their finite-case
elimination. They do not become additional permissible $F$ bodies.
The compiler of the preceding section is used whenever an iteration
update refers to one of these graphs.

\subsection{Difference and doubling}
Use Bounded iteration with initializer $x$, update $P$, and cap $x$.
Since $P(v)\le v$, clipping is never active. This defines a total
finite-case graph $d(x,n)$ with
\[
d(x,0)=x,\qquad d(x,Sn)=P(d(x,n)),\qquad d(x,n)\le x.
\]
Equation induction on $y$, with $x$ fixed, proves
$d(Sx,Sy)=d(x,y)$: its base is $P(Sx)=x$, and its step applies $P$
to the previous equality. In particular $d(n,n)=0$ by a further
equation induction. No interpretation as subtraction is assumed
at this stage.

Let
\[
 Q(n):=F(n;v.S^2v),\qquad \operatorname{db}(n):=P^3(Q(n)).
\]
Here $Q$ is a source term, and $\operatorname{db}$ is a finite-case
graph. Their equations are
\[
 Q(0)=3,\quad Q(Sn)=S^2Q(n),\qquad
 \operatorname{db}(0)=0,\quad
 \operatorname{db}(Sn)=S^2\operatorname{db}(n),\quad
 Q(n)=S^3\operatorname{db}(n).
\]
For the last three identities, first prove $Q(n)\ge3$ by induction;
functional predecessor coverage then permits cancelling precisely
three successors, and the displayed recurrence follows. Quantifier-free
induction proves $n\le\operatorname{db}(n)$ and monotonicity. For
monotonicity induct on the larger argument in
$a\le b\to\operatorname{db}(a)\le\operatorname{db}(b)$ and use
the equality-or-predecessor alternatives in New-cap.

\subsection{Addition by a saturated diagonal argument}
For any cap $c$, let $A_c(n,a)$ be Bounded iteration with initializer
$a$ and update $S$. Write $s_c(v)=\min(Sv,c)$. Thus
\[
 A_c(0,a)=\min(a,c),\qquad A_c(Sn,a)=s_c(A_c(n,a)).
\]
For $a\le c$, equation induction gives seed commutation:
\[
                  A_c(n,s_c(a))=s_c(A_c(n,a)).
 \tag{Seed commute}
\]
The two sides have the same base and the same update. The finite-case
minimum graph obeys $\min(0,c)=0$ and
$\min(Sn,c)=s_c(\min(n,c))$, by the discrete linear-order laws.
Iteration uniqueness therefore gives
$A_c(n,0)=\min(n,c)$.

Fix $a\le c$ and induct on $n$ in
\[
 n\le c\ \longrightarrow\ A_c(n,a)=A_c(a,n).
\]
The base follows from the preceding minimum identity. Whenever $Sn\le c$,
the step applies $s_c$ on the first side, and Seed commute on the
second; the induction hypothesis then identifies them. This proves
commutativity for bounded arguments without assuming addition.
The increasing update and order-preserving $s_c$ also prove that
$A_c$ is monotone in its counter and in its initializer within $[0,c]$.
Counter monotonicity uses induction on the larger counter and New-cap;
initializer monotonicity is direct induction on the common counter.

Now fix an input bound $b$ and choose the \emph{source} cap
$c=Q(b)=S^3\operatorname{db}(b)$. Quantifier-free induction on $n$ in
\[
 n\le b\ \longrightarrow\ A_c(n,n)=\operatorname{db}(n)
\]
proves the diagonal identity. At a successor Seed commute makes
$A_c(Sn,Sn)=s_c^2(A_c(n,n))$. Since
$\operatorname{db}(Sn)\le\operatorname{db}(b)<c$, neither of these
two successor operations is clipped. Thus the successor diagonal
equals $S^2\operatorname{db}(n)=\operatorname{db}(Sn)$.

Monotonicity now gives the global local bound
\[
 n,a\le b\quad\Longrightarrow\quad
 A_{Q(b)}(n,a)\le A_{Q(b)}(b,b)
   =\operatorname{db}(b)<Q(b).
 \tag{Addition bound}
\]
So these runs have no clipping on bounded inputs. For two caps from
input bounds $b,b'$, induction on $n$ restricted by $n\le b,b'$
proves that the results agree whenever $a\le b,b'$: the two bases
are $a$, and both successor steps are genuine successor operations.

Choose $b$ as the first largest of the two arguments, by the finitely
many source order cases. This defines a finite-case graph $n+a$.
Coherence proves its defining equations
$0+a=a$ and $(Sn)+a=S(n+a)$ across possible changes of chosen cap.
The graph is commutative by bounded commutativity. It has the
source unary bound $Q(b)$ on arguments at most $b$.
Usual associativity and monotonicity follow by fixed-parameter
quantifier-free inductions from these equations.

\subsection{Multiplication and primitive-recursive pairing}
Define the source unary function
\[
 E(n):=F(n;v.Q(v)).
\]
It satisfies $E0=4$, $E(Sn)=Q(E(n))=\operatorname{db}(E(n))+3$,
is monotone, and $E(n)\ge n$. Fix $n,a\le b$ and apply Bounded
iteration with initial value zero, update $v\mapsto v+a$, and
source cap $E(Q(b))$. For the run at stage $j\le b$, use the
quantifier-free invariant
\[
                         v_j\le E(b+j).
\]
Its base is immediate. At a successor,
\[
 v_j+a\le E(b+j)+b
          \le\operatorname{db}(E(b+j))<E(S(b+j)).
\]
Here $E(b+j)\ge b$ and the already constructed addition is used
only through its finite-case graph. Also
$b+j\le\operatorname{db}(b)<Q(b)$, so the stage bound is at most
the chosen cap. Thus clipping is inactive throughout the intended run.
The same cap-coherence induction as for addition gives a finite-case
multiplication graph, with source unary bound
\[
                         M_\times(b)=E(Q(b)).
\]
Its zero and successor equations give the elementary nonnegative
semiring identities by quantifier-free induction. In particular,
$w^2\ge w$ for $w\ge1$ and
$w^2+w<(Sw)^2$.

Put
\[
                         \kappa(x,y)=(x+y)^2+x.
\]
This is a finite-case graph, with a monotone inflationary source bound
\[
                  K(b)=Q(M_\times(Q(b))).
\]
Both coordinates of a valid code are no greater than the code.
To construct its projections, use capped scalar iteration starting
at zero with update
\[
 u\longmapsto\begin{cases}Su,&(Su)^2\le z,\\u,&\text{otherwise},\end{cases}
\]
cap $z$, and counter $z$. The update preserves $u\le z$ and
$u^2\le z$. The additional quantifier-free invariant
\[
               u_j=j\quad\lor\quad\text{the update fixes }u_j
\]
shows that the state at $j=z$ is fixed: if it equals $z$, then
$(Sz)^2>z$ already makes it fixed. Hence the result $r(z)$ satisfies
$r(z)^2\le z<(Sr(z))^2$. Linear order and square monotonicity make
such a result unique.

Set
\[
 \pi_0(z)=d(z,r(z)^2),\qquad \pi_1(z)=d(r(z),\pi_0(z)).
\]
The difference identity $d(x+y,y)=x$ follows by equation induction
on $y$, using $d(Sx,Sy)=d(x,y)$ and the addition equations.
For $z=\kappa(x,y)$, put $w=x+y$. The inequalities
$w^2\le z\le w^2+w<(Sw)^2$ give $r(z)=w$, whence
\[
               \pi_0(\kappa(x,y))=x,\qquad
               \pi_1(\kappa(x,y))=y.
\]
The graph functions are total on arbitrary codes too, and
$\pi_0(z),\pi_1(z)\le z$, since each difference is bounded by
its first argument. No unbounded minimization has been introduced.

\section{Full primitive recursion and both interpretations}
\label{sec:selfseed-equivalence}

\subsection{A source bound for every fixed description}
We externally recurse on ordinary finite primitive-recursive
descriptions. The induction data for $f$ are a finite-case graph,
internal proofs of its defining equations, and a monotone inflationary
\emph{source unary term} $M_f$ such that
\[
              \vec x\le b\quad\Longrightarrow\quad f(\vec x)\le M_f(b).
\]
Zero, successor, and projections have these data with the bound $Q$.
Composition takes finite common refinements of cases. If
$f=h(g_1,\ldots,g_m)$, the composition
$M_h\circ M_{g_1}\circ\cdots\circ M_{g_m}$ is a suitable source
bound: monotonicity and inflationarity make it dominate each needed
intermediate argument. Empty tuples and nullary descriptions are
handled by their fixed finite source values and inflationary constant
shifts of $b$.

Suppose the next description is
\[
f(0,\vec p)=g(\vec p),\qquad
f(Sn,\vec p)=h(n,f(n,\vec p),\vec p).
\]
Let $G=M_g$, $H=M_h$, and choose the source unary function
\[
D=K\circ H\circ G\circ Q,\qquad
U(t)=F(t;v.D(v)),\qquad B(b)=U(SQ(b)).
\]
Every factor is monotone and inflationary, and $D(v)\ge Sv$.
Consequently $U$ is monotone and $U(t)\ge St$; $B$ is a monotone
inflationary source term. These facts are quantifier-free inductions
using the already proved order laws.

The scalar state update is the finite-case graph
\[
 J(z,\vec p)=\kappa\bigl(S\pi_0(z),
                    h(\pi_0(z),\pi_1(z),\vec p)\bigr),
\]
with initializer $\kappa(0,g(\vec p))$. Apply the bounded source
compiler to this initializer and update with cap $B(b)$. Assume
$n,\vec p\le b$. For stages $j\le b$ prove
\[
                          z_j\le U(Sb+j).
 \tag{PR-stage}
\]
The seed satisfies
\[
\kappa(0,g(\vec p))\le K(G(b))\le D(b)
                         \le D(U(b))=U(Sb).
\]
For the successor, write $s=U(Sb+j)$. Then $s\ge b$, and both
state projections are at most $z_j\le s$. The next counter is
at most $Ss\le Q(s)$, and the new value is at most
$H(G(Q(s)))$, by monotonicity and inflationarity of the bounds.
The pairing bound therefore gives
\[
                   J(z_j,\vec p)\le D(s)=U(S(Sb+j)).
\]
Since $j\le b$, the stage argument $Sb+j$ is at most
$S\operatorname{db}(b)$ and hence at most $SQ(b)$. Thus every
stage bound is at most $B(b)$, proving that no clipping occurs.
The displayed invariant is quantifier free in the finite-case
definitional extension; it translates to permitted source induction.

The projection identities prove by quantifier-free induction guarded
by $j\le b$ that the counter component of $z_j$ is $j$, and that the
value component satisfies its required base and, when $Sj\le b$,
step equations. Two valid caps
give the same run up to their common allowed counter by induction:
the seeds agree and their unclipped updates are identical. Choose
$b$ as the first largest input among $n,\vec p$, by finitely many
source order cases, and define $f(n,\vec p)=\pi_1(z_n)$. This is
again a finite-case graph. At any larger input bound $b$, cap coherence
identifies this value with the run at cap $B(b)$, or monotonicity of
$B$ bounds it from the selected smaller input bound. Thus its value
is bounded by $B(b)$ for every valid input bound, and take $M_f=B$.
The cap-coherence argument proves its defining equations
even when a successor input changes the selected largest coordinate.

The source compiler has expanded every use of $g,h,\kappa,\pi_i$
inside the scalar update into original-language terms, encoded cases,
and bounded selectors. Neither a total evaluator for description
codes nor a new recursor schema over graph-defined operations has
been assumed. The construction is a finite external recursion on
each fixed description.

\subsection{Whole-domain bi-interpretation}
Interpret the source theory in iPRA on its full numerical domain.
Keep $0,S$ and actual equality; take $\epsilon$ to be the standard
primitive-recursive equality characteristic. For each fixed earlier
body $h$, interpret $F_h$ by primitive recursion with seed $S h(0)$
and scalar update $h$. All source axioms and its equation-induction
rule then hold in iPRA.

Conversely, interpret each fixed PR function description by the
finite-case graph constructed above. Its totality, single-valuedness,
defining equations, and quantifier-free induction instances have all
been proved in the source. The original successor, injectivity,
separation, and actual equality are retained. Thus this is an
interpretation of full first-order iPRA, with its ordinary intuitionistic
logic and the entire numerical domain.

For the iPRA-side composite, external induction on a fixed PR
description reduces identification with the original function to
uniqueness from its base and step equations. Internal quantifier-free
induction proves this uniqueness at arbitrary counters and parameters.
For the source-side composite, $0,S$ are unchanged. Both equality
characteristics have the two-value range and reflect actual equality,
so Boolean cases identify them. For a fixed $F_h$, external induction
on its finite stage identifies the earlier body; both the old and
composite functions then satisfy the same self-seeded base and scalar
step equations. Equation induction identifies them everywhere. If a
composite uses graph-defined functions, the finite-case elimination
lemma supplies the required source induction instance.

The relation $x=y$ therefore defines the isomorphism of each domain
composite with the original structure. External induction on formula
syntax extends preservation of functions and equality to all formulas;
quantifiers range over the same whole domain and genuine falsity is
preserved. This gives whole-domain bi-interpretability and
equiconsistency. It is not an internal assertion of either theory's
own consistency.

\paragraph{Research status and complexity.}
This is a written proof of a new scalar self-seeded basis, with the
separately recorded independent mathematical audit. It removes both
the independent initializer and the old counter argument from the
primitive recursor. It pays explicitly for actual numerical equality,
successor injectivity, and ordinary intuitionistic
logic. No complete checker for this new theory has been measured,
and no existing tiny initialized-iterator checker is assigned to it.
It must not be described as a demonstrated improvement below the
current complete-checker frontier or as an absolutely simplest theory.
